\pagebreak

\chapter{Constraints, Properties and Attributes}\label{chap:cpa}
    \index{constraint}
    \index{property}
    \index{attribute}
    
    These three terms are used somewhat interchangeably by Xilinx.
    Some constraints/properties/attributes are generated automatically by 3PL.
    These are -
    \begin{enumerate}
    \item   input and output port locations
    \item   CLB memory initialisation
    \item   block memory initialisation
    \end{enumerate}
    These are written to the EDIF file as properties on instances of elements.
    
    All other constraints are generated by 3PL code, in the case of Xilinx in
    library xilinx/common.3pl. The constraints are written to the NCF file in the
    case of ISE or an XDC file for Vivado.
    
    Input and output connections have location constraints generated if a {\bf
    loc} attribute has been set for the input or output mode variable. Attributes
    {\bf ibufdelay}, {\bf ifddelay}, {\bf tnm}, {\bf pullup}, {\bf pulldown},
    {\bf differential}, {\bf diff\_term}, {\bf iostandard}, {\bf ibufg}, {\bf
    slew}, {\bf drive}, if set, also result in constraint file entries.  
    Inbuilt procedures {\bf input()} \autorefph{sec:ipinput}, {\bf
    output()} \autorefph{sec:ipoutput} and {\bf clock()} \autorefph{sec:ipclock} do
    not directly generate these constraints; for convenience and flexibility they
    are handled in 3PL software.

    While attributes can be passed directly to inbuilt procedures {\bf input()},
    {\bf output()} and {\bf clock()}, they may also be associated with the input pins. This
    is done by using the library procedure {\bf iopins()}
    \autorefph{sec:lpiopins} to create variables which contain strings naming the
    FPGA I/O pins (locations). {\bf iopins()} also creates a global map {\bf
    pinAttributes} whose entry keys are I/O pin names, the associated values
    being attribute maps whose keys are attribute identifiers and values the
    attribute values.

    A 3PL user-defined procedure to extract these pin-associated attributes can be
    automatically called whenever an input, output or clock mode variable is
    created. The Xilinx 3PL libraries contain such a procedure, {\bf
    processPinAttributes()} \autorefph{sec:lpprocPinAtt}. This procedure reads the
    pin locations (loc attribute) from the variable and then reads the attribute
    map for each pin, adding these attributes to the attribute map for the
    variable. Note that these pin-related attributes will be overridden by
    attributes passed as arguments to input(), output() or clock(). {\bf
    processPinAttributes()} also ensures that attributes are converted from single
    value to per-pin arrays where appropriate.

    The above procedure must be registered with 3PL so that it can be called by
    the system when an input, output or clock mode variable is created. This is
    done by calling inbuilt procedure {\bf attributeprocedure()}
    \autorefph{sec:ipattributeproc} with a pointer to the procedure as
    argument. Clearly this should be done before any of the above I/O procedures
    are called, preferably be placing a call to {\bf attributeprocedure()} early
    in a leader module.

    While netlist generation attributes are handled by input(), output() of
    clock() following completion of immediate execution, constraint file
    attributes must be handled by 3PL software. This is done by library procedure
    {\bf generateConstraints()} \autorefph{sec:lpgenerateconstraints}. This
    procedure should be called late in a trailer module to ensure that no more
    calls to input(), output() or clock() will be executed. This procedure calls
    inbuilt function {\bf varlist()} \autorefph{sec:ifvarlist} which returns a
    list of target mode variables. The resulting list is traversed and the
    constraint-related attributes of input, output and clock mode variables are
    extracted and constraints are written to the constraints file.   
    
    An example of constraints that might need to be generated explicitly are
    input and output delays where combinatorial logic paths are included - see
    \autorefph{sec:input_output}.
    
    The clock timing group identifier can be explicitly set using a clock
    attribute (tnm), otherwise it will be constructed from the
    clock identifier. Regardless of whether a frequency or period attribute
    has been set a period constraint will be generated. If no frequency
    or period attribute has been set a warning message will be produced to
    the effect that there will be no timing constraints for that clock.
    
    CLB RAM and block RAM initialisation entries are added to the
    EDIF file. If no initialisation is supplied to CMEMORY or RMEMORY
    mode variables the initialisation entries will be 0.

    The library {\bf common.3pl} contains procedures for adding constraints to the
    NCF and XDC file.

    The following raw low level procedures can be used for general constraint strings
    or constraint strings applying to netlists, but for most constraints there are
    more convenient ways to apply them - \\
    {\bf addncfline()} \autorefph{sec:lpaddncfline} \\
    {\bf netaddncf()} \autorefph{sec:lpnetaddncf} \\
    {\bf elementaddncf()} \autorefph{sec:lpelementaddncf} \\
    {\bf addxdcline()} \autorefph{sec:lpaddxdcline}
    
    I/O constraints such as slew rate, drive current or input delays are normally
    passed to inbuilt procedures {\bf input()} \autorefph{sec:ipinput} and {\bf
    output()} \autorefph{sec:ipoutput} as attributes from which the constraint file
    entries are automatically generated. If I/O logic is being built from the bottom
    up using inbuilt procedure {\bf element()} and associated procedures then I/O
    constraints can be applied by library procedure {\bf IOconstraint()}.
    \autorefph{sec:lpioconstraint}.

    
    Direct creation of elements is
    usually done in library procedures, for example in xilinx/clocks.3pl,
    but this can also be done in any application code using inbuilt procedure
    {\bf element()} \autorefph{sec:ipelement}.
    Where a vendor library element has been explicitly created by 3PL code
    properties can be added by using inbuilt procedure 
    \begin{lstlisting}[gobble=4, language=threepl]
       elementproperty(id, property, value);
    \end{lstlisting}
    where all three arguments are strings. There are more convenient library
    procedures to call this indirectly - \\
    {\bf elementLogProperty()} \autorefph{sec:lpelementLogProperty}, \\
    {\bf elementIntProperty()} \autorefph{sec:lpelementIntProperty}, \\
    {\bf elementBinProperty()} \autorefph{sec:lpelementBinProperty}, \\
    {\bf elementHexProperty()} \autorefph{sec:lpelementHexProperty}, \\
    {\bf elementFloatProperty()} \autorefph{sec:lpelementFloatProperty}, \\
    {\bf elementStrProperty()} \autorefph{sec:lpelementStrProperty}.

    The main netlist constraints are
    \blist
    \item    I/O pad location
    \item    absolute element location constraints
    \item    relative element location constraints
    \item    time constraints
    \blend

    Absolute location constraints are not inherently supported by 3PL except for I/O
    pads and their associated buffers and I/O registers (static mode variables).
    Single objects such as global clock buffers or clock managers can be given
    absolute locations by calling inbuilt procedure {\bf elementproperty()} with a
    "LOC" property, or via library procedure {\bf elementStrProperty()}. registers
    could be given location constraints by reading back block ID attributes from the
    static mode variable and applying location constraints to these.
    
    Relative location constraints are not explicitly supported at the moment but
    they could be generated by application code using constraint file
    procedures described above.
    
    3PL takes high-level statements and synthesises logic to implement these
    statements using basic logic elements. Since the logic synthesis process
    is not accessible to the programmer there is no way to supply location
    information pertinent to synthesised logic elements. In extreme cases
    logic can be constructed using inbuilt procedure {\bf element()} and
    location constraints can be attached to these, however this is very low level
    logic construction.
    
    Clock managers (dll, dcm, pll, clkmmcm) are created by library procedures which
    can be passed a location argument.
    
    Global clock buffers can be created by using library procedure {\bf
    clockbufg()} in {\bf xilinx/clocks.3pl} (all procedures in this library
    have optional location arguments). This library procedure is an example of the
    use of procedures {\bf element()}, {\bf elementinput()} and {\bf elementoutput()} and its code is
    an example of the use of these procedures -

    \begin{lstlisting}[gobble=4, language=threepl]
       global procedure
       clockbufg (clock(out) <- clock(in), str(loc), log(invert)) {
           elementCheck("BUFG", "clockbufg");
           element("BUFG");
           if (invert) {
               value(_in, "log", !in);
               elementinput(, "I", _in);
           } else
               elementinput(, "I", in);
           elementoutput(out <- , "O");

           attributes(out,
               frequency=frequency(in),
               period=period(in),
               dutycycle=dutycycle(in)
           );
           if (matched(loc)) elementStrProperty(, "LOC", loc);

           elementend();
       }
    \end{lstlisting}
    The first input argument is omitted in each call to elementinput()  and
    elementoutput() since we have not bothered using an element identifier.    
   
    Input arguments to parameters {\bf loc} and {\bf invert} are optional. The
    procedure calls {\bf element()} to create an element called {\bf BUFG} (a
    Xilinx global clock buffer). If {\bf invert} is false (which it is by default
    if not passed an argument) then input argument {\bf in}, a clock signal, is
    connected to pin {\bf I} port type 0 (input). If {\bf invert} is true an
    inverted signal {\bf \_in} is created (value mode variables of type {\bf log}
    can be used with clock mode variables) which inverts {\bf in} and is connected
    to the buffer. Likewise parameter {\bf out} is connected to pin
    {\bf O} port type 1 (output).
    
    The next statement applies an absolute location constraint to the clock buffer
    if parameter {\bf loc} has been supplied an argument. For example the argument
    might be the string "GCLKBUF2" in which case library procedure
    {\bf elementStrProperty()} is called with arguments "LOC" and "GCLKBUF2". This
    library procedure simply makes the inbuilt procedure call
    {\bf special(0, "LOC=GCLKBUF2")} which prepends "INST " and the element
    identifier to the location string and writes it to the NCF file.
    
    The last three statement copies the frequency, period and duty cycle of the
    input parameter clock variable to the output parameter clock variable.
        
    Location constraints can be applied to elements explicitly instantiated by
    using inbuilt procedures {\bf element()}. When using {\bf element()} the
    example above shows the adding of a location constraint by using library
    procedure {\bf elementStrProperty()}.
    
    As a further example suppose using ISE we wish to register an input signal from pad
    "A1" in a flip-flop at location "SLICE\_X12Y10" (I don't believe it is possible to
    specify the X or Y flip-flop in the slice!).
    \begin{lstlisting}[gobble=4, language=threepl]
       input(in, "log", loc="A1");
       value(out, "log");
       element("fdre");
       elementclockinput(, "C", c100); // connect clock
       elementinput(, "D", in);        // connect data in
       elementoutput(out <- , "Q");    // connect data out
       elementaddncf("LOC=SLICE_X12Y10");
       elementend();
    \end{lstlisting}
    The value mode variable {\bf out} can now be used in 3PL as usual. The flip-flop
    will be clocked on every clock cycle, but a clock enable input could
    be connected.
    
    To do the same as above using Vivado the location constraint would have to be
    generated differently, for which an explicit element identifier would be needed -
    \begin{lstlisting}[gobble=4, language=threepl]
       input(in, "log", loc="A1");
       value(out, "log");
       element("fdre", "dff123");
       elementclockinput(, "C", c100); // connect clock
       elementinput(, "D", in);        // connect data in
       elementoutput(out <- , "Q");    // connect data out
       elementend();
       addxdcline("set_property LOC SLICE_X12Y10 [get_cells dff123]");   
    \end{lstlisting}    
    
    In the above example if we did not have the location constraint for the
    flip-flop we could have written it in 3PL as -
    \begin{lstlisting}[gobble=4, language=threepl]
       input(in, "log", loc="A1");
       static(s, "log");
       useclock(c100);
       s <- in;
    \end{lstlisting}
    However, there is one important distinction - the high-level version will
    not start executing until it gets a start signal (this may be from an
    external source or may be automatic after configuration loading) as it is
    executing in a thread within an infinite loop initiated formally within 3PL.
    The previous low-level example will start executing immediately the clock is
    available.    
    
    \index{constraint!timing}
    \index{timing constraint}
    Time constraints are already handled well in 3PL. These are usually set
    up in platform header files so that application programs can ignore
    timing constraints.
    
    When clock variables are created they have a frequency or period specified.
    The clock period is used to automatically provide a timing constraint for
    all logic driven by that clock (inbuilt procedure {\bf useclock()} sets the
    current clock domain for all following logic). However, there are some
    situations in which timing constraints are not automatically generated -
    \blist
    \item   timing constraints between I/O pads and clocked elements
    \item   timing constraints between clock domains
    \item   situations where a timing constraint can be longer than that of
            the period of the clock in use
    \blend
    
    Timing constraints between I/O pads and state elements are illustrated
    in \autorefph{sec:input_output}.

    Timing constraints between clock domains is only meaningfull if the clocks
    are related in frequency and have a known phase relationship. 3PL normally
    forbids data from one clock domain being being simply assigned in another,
    but it provides three mechanisms to avoid this. Firstly queues can be read
    and written in different clock domains. The logic implementing this uses
    asynchronous FIFOs and there need be no relationship between the two clock
    domains. Secondly the inbuilt function {\bf sample()} allows data in one
    clock domain to be sampled in another. Again there need be no relationship
    between the two clock domains. Thirdly the inbuilt function {\bf accept()}
    can be used to force simple assignment to be accepted by the compiler. This
    can only be used safely under certain circumstances when the clock domains
    are related and have a known relative phase. In this case a timing
    constraint must be used between state elements in the two clock domains. A
    common example is where one clock is a multiple of the other and they are in
    phase, i.e. the positive edge of the slower clock coincides with a positive
    edge of the faster clock. For example if we have two clocks {\bf c100} and
    {\bf c50} which are 100MHz and 50MHz respectively and are in phase we can
    assign from one domain into the other -
    \begin{lstlisting}[gobble=4, language=threepl]
       maxdelay("C100", "C50", 10.0e-9);
       static({v1, v2}, type);
       useclock(c100);
       v1 <- .....;        // as side effect sets the clock domain of v1 to c100
       useclock(c50);
       v2 <- accept(v1);   // assigning v1 in domain c50 - accept() prevents fatal error
    \end{lstlisting}
    
    When a clock variable is created it will be given a Xilinx timing group
    identifier which by default is the variable identifier converted to upper
    case. During place and route all state elements clocked by a clock are
    entered in a set associated with the clock timing group identifier. Data
    paths between state elements in the same timing group are automatically
    constrained to be no greater than the clock period minus setup time. Date
    paths between state elements in different timing groups are not constrained
    by default. The example above applies a timing constraint along data paths
    from state elements in timing group {\bf C100} to state elements in timing
    group {\bf C50}. The constraint is 10.0ns and for reporting purposes the
    constraint will be given the identifier "TS\_C100\_C50".

    The application of arbitrary timing constraints between registers, for
    example where a register might be assigned every second clock cycle, is
    probably not a good idea. This is not in the spirit of the language and
    would involve code which was difficult to verify and maintain. It can be
    done by reading the block identifier attributes of the registers concerned
    and applying a timing constraint between them.
